\section{Installation}
\label{sec:installation}

\denis{} runs on Windows, macOS, and Linux. The recommended install
flow uses \href{https://docs.astral.sh/uv/}{\code{uv}} (a fast Python
package and project manager from Astral) to manage Python itself, the
project virtual environment, and all dependencies in one step. No
existing Python installation is required: \code{uv} downloads a
compatible interpreter on first sync.

\subsection{Recommended: scripted install}

After cloning the repository, open a terminal in the project root and
run the platform-specific install script.

\begin{lstlisting}
# Windows (cmd or PowerShell)
install.bat

# Linux / macOS
./install.sh
\end{lstlisting}

The script:
\begin{enumerate}
    \item Checks whether \code{uv} is on \code{PATH}; if not, prompts to
        install it via the official Astral installer.
    \item Runs \code{uv sync}, which downloads a compatible Python
        interpreter (if needed), creates \code{.venv/} in the project
        root, and installs every dependency pinned in
        \code{uv.lock}.
    \item Creates a desktop shortcut (also available on its own via
        \code{install.bat shortcut} / \code{./install.sh shortcut}):
        a Windows \code{.lnk} pointing at
        \code{.venv\textbackslash Scripts\textbackslash pythonw.exe},
        a Linux \code{.desktop} entry, or a clickable
        \code{.command} on macOS.
\end{enumerate}

When the script finishes, double-click the desktop shortcut
(Figure~\ref{fig:installation_shortcut}) to launch \denis{}, or run
from the terminal:

\begin{lstlisting}
uv run python gui.py
\end{lstlisting}

The first launch opens a splash screen, then the main window with the
four tabs introduced in Section~\ref{sec:introduction}.

\begin{figure}[hbtp]
    \centering
    \includegraphics[width=2.4cm]{fig_installation_shortcut.png}
    \caption{The \code{DENIS} desktop shortcut created on Windows by
        the installer. The Linux \code{.desktop} entry
        and macOS \code{.command} launcher use the same icon.}
    \label{fig:installation_shortcut}
\end{figure}

\subsection{Manual install}

If \code{uv} is already installed and the wrapper scripts are not
needed:

\begin{lstlisting}
git clone https://github.com/ardakayaalp/DENIS.git
cd DENIS
uv sync
uv run python gui.py
\end{lstlisting}

The \code{clstools}~\cite{clstools} dependency is not published on
PyPI and is fetched from GitHub via the \code{[tool.uv.sources]}
block in \code{pyproject.toml}. \code{uv sync} handles this
automatically; no manual step is needed.

\subsection{Editor / IDE integration}

For VSCode or PyCharm, point the project Python interpreter at the
\code{uv}-managed environment so debugging and linting use the same
packages as the running app:

\begin{itemize}
    \item Windows: \code{.venv\textbackslash Scripts\textbackslash python.exe}
    \item Linux / macOS: \code{.venv/bin/python}
\end{itemize}

\subsection{Updating}

To pull a newer version of \denis{} and refresh dependencies:

\begin{lstlisting}
git pull
uv sync
\end{lstlisting}

\code{uv sync} reads \code{uv.lock} and brings the local environment
to exactly that pinned set --- adding, removing, and upgrading
packages as required.

\subsection{Troubleshooting}

\begin{description}
    \item[\code{uv} installed but not found.]
        The Astral installer drops \code{uv} into
        \code{\textasciitilde/.local/bin} (Linux/macOS) or
        \code{\%USERPROFILE\%\textbackslash .local\textbackslash bin}
        (Windows). Open a new terminal so the updated \code{PATH}
        takes effect, or add the directory to \code{PATH} manually.
    \item[\code{clstools} fetch fails.]
        \code{uv sync} clones \code{cls\_tools} from GitHub. A
        failure usually means the machine has no GitHub access. Clone
        the repository to a sibling folder manually and replace the
        entry in \code{[tool.uv.sources]} with a local path:
\begin{lstlisting}
[tool.uv.sources]
clstools = { path = "../cls_tools" }
\end{lstlisting}
    \item[Desktop shortcut opens a console on Windows.]
        \code{create\_shortcut.bat} prefers
        \code{.venv\textbackslash Scripts\textbackslash pythonw.exe};
        if it falls back to \code{python.exe}, \code{pythonw.exe} is
        missing from the environment. Re-run \code{uv sync} to repair.
    \item[macOS \code{.command} launcher will not open.]
        macOS may quarantine downloaded scripts. Right-click the
        \code{DENIS.command} file, choose \app{Open}, and allow
        execution at the Gatekeeper prompt; or grant execute
        permission with \code{chmod +x \textasciitilde/Desktop/DENIS.command}.
\end{description}
